\section{결론}\label{sec:conclusion}
본 연구는 영상 기반 초기 추정기를 고정한 채, 다중 해상도 특징과 알려진 물리적 치수로 코너의 국소 이동을 학습하는 보정 방법을 평가하였다. 이동 후보의 확률 가중평균과 원본 좌표의 상한을 사용하고, 동일한 자세 계산을 통해 코너 수정과 위치·회전 변화의 관계를 비교하였다.

직사각형 실사 319장에서 세 기반의 코너 중앙 오차와 전체 주석 기준 \pckten 이 개선되었으며, 별도 정사각형 119장에서도 2D 보정 효과가 관찰되었다. 반면 ResNet 위치 중앙값의 사후 구간은 0을 포함했고, 일부 P90과 사람 입력의 어려움 집단의 위치 오차는 악화되었다. 회전 대응 감독의 추가 효과도 일관되지 않았다. 네 기록 영상 8,910프레임에서는 Base와 N3의 자세 산출률이 같은 98.451\%였고 인접 출력 변화는 지표별로 혼합되었다. 이러한 결과는 주어진 정보·학습·참조 조건에서 국소 보정의 적용 가능성과 경계를 함께 보여준다. 정사각형을 포함한 사람 입력 등급과 3,101개 참조 코너의 가시성 상태는 확보하였다. 별도 12장 PnP 보조 기하 참조에서는 P90과 PCK가 개선되었지만 중앙값은 악화되어 결과가 혼합되었다. 이 탐색 패널의 실제 대상 대응 12건은 확보했으나 공식 120/24장 가시 코너 계약과 독립 물리 자세 정확도를 대신하지 않는다. 남은 확인은 등급 판정 기준·참조 품질과 독립 물리 참조이며, 기록 영상의 정지 잡음도 확인된 정지 구간이 필요하다. 이미 입력된 등급과 가시성의 재작업을 전제로 하지 않으며, 추가 일반화는 독립 관측에서 평가해야 한다.

추가 모델 변경의 필요성과 독립 측정의 필요성은 분리된다. 기존 N3 결과와 대안의 중앙값·꼬리·실패·비용은 현재 원고에서 정리할 수 있지만, 반복 개발자료의 수치만으로 독립 물리 정확도나 실제 가림 대응을 확정할 수 없다. 제한된 자세 결합 후보 실험의 6회 학습·세 seed 비교·동일 경계 비용도 보충자료에 보존하였다. 후보 oracle 여지가 있어도 새 학습은 실용 N3/PoseFix 대비 개선을 지지하지 않았으며, 새로운 확증에는 같은 상대 자세의 실물 가림 쌍과 독립 계측 참조가 필요하다.
